\section{Background}
\label{sec:background}

Estimating how frequently a rare event recurs is a classical problem in
applied statistics, arising in contexts ranging from seismology to
equipment failure monitoring. The canonical model treats event occurrences
as realizations of a point process on the real line, and the recurrence
interval is characterized by the distribution of inter-arrival times under
that process \citep{cox1966}. When the process is assumed to be a
homogeneous Poisson process, the inter-arrival times are exponentially
distributed and the recurrence interval reduces to the reciprocal of a
single rate parameter, which can be estimated directly from the observed
event counts. In practice, however, sensor networks deployed for rare-event
monitoring rarely satisfy the homogeneity assumption: sensitivity drifts
over time, environmental conditions modulate the true event rate, and
individual sensors intermittently drop out of the network, all of which
violate the assumptions underlying the simplest estimators.

\subsection{Point Process Models}
A more flexible family of models represents the event rate as a
time-varying intensity function $\lambda(t)$, so that the expected number
of events in an interval $[t_1, t_2]$ is given by the integral of
$\lambda(t)$ over that interval, following the general treatment
of~\citet{daley2003}. Self-exciting variants of this model, most notably
the Hawkes process, further allow each observed event to temporarily elevate
the intensity of future events, which is a useful property when modeling
phenomena such as aftershock sequences or cascading equipment failures.
Bayesian inference procedures for Hawkes processes have been studied
extensively; see \citet{rasmussen2013} for a thorough treatment of
posterior sampling schemes under weakly informative priors.\footnote{In our own experiments we found that a weakly informative
gamma prior on the baseline intensity, combined with a Beta prior on the
excitation decay rate, gave stable posterior estimates even with as few as
twenty observed events.}

\subsection{Sparse Sensor Networks}
The setting considered in this thesis differs from the classical point
process literature in one important respect: observations do not arrive
from a single continuously monitored location, but rather from a spatially
distributed network of low-power sensors that report intermittently and
that individually observe only a fraction of the true event population.
\citet{fictional2020} studied the related problem of sensor placement for
rare-event monitoring and showed that even a modest number of well-placed
sensors can recover an unbiased estimate of the aggregate event rate,
provided that the missingness process is independent of the underlying
event intensity. This independence assumption is central to the estimator
developed in Section~\ref{sec:method}, and we return to it explicitly when
discussing the conditions under which our estimator remains consistent.

Figure~\ref{fig:network} illustrates a representative deployment topology
used in the field study described later in this chapter, showing sensor
nodes placed at irregular intervals along a monitored corridor.

\begin{figure}[h]
\centering
\includegraphics[draft,width=0.7\linewidth]{network.pdf}
\caption{Representative sparse sensor deployment topology used in the
field study. Circles denote sensor nodes; the shaded region denotes the
monitored corridor.}
\label{fig:network}
\end{figure}

As the figure suggests, the effective sampling density varies considerably
across the monitored region, which motivates the spatially adaptive
weighting scheme introduced later in the chapter. The remainder of this
chapter builds toward a single unified estimator that accounts for both the
temporal irregularity of event arrivals and the spatial sparsity of the
sensor network, combining ideas from the classical point-process literature
with the placement-aware corrections proposed by \citet{fictional2020}.
